\section{Conclusion}
\label{sec:conclusion}

% Included in both manuscript versions on 2026-11-03 at the author's request.
In this paper, we investigated GNSS-independent GFRA for LEO satellite
IoT systems. The proposed multi-frame AFDM scheme
accommodates large differential delays and Doppler shifts without
device-side precompensation, while its compact effective-channel
representation allows observations from additional pilot frames to be
combined without increasing the number of unknown channel coefficients.
Building on this representation, we developed the GP-MP receiver to
jointly infer device activity and effective channels. Channel estimation
exploits a score-based generative prior learned from effective-channel
samples without imposing sparse or grid-aligned assumptions. The
analytical relationship between the posterior mean and variance further
allows the same denoiser to provide the channel estimates and uncertainty
information required by the receiver, without a separately trained
variance model. These estimates are combined with activity probabilities
obtained through density-ratio estimation, which evaluates the activity likelihood
ratio without a Gaussian approximation.
Simulation results demonstrated that this transmission and receiver
design improves both activity detection and channel estimation over the
adapted STMP baseline. These gains also reduce the required number
of pilot frames: with fixed energy per frame, GP-MP reached the
$-12$~dB NMSE target with 16 pilot frames, compared with 20 for STMP.
